% Part: computability
% Chapter: computability-theory
% Section: fixed-point-thm

\documentclass[../../../include/open-logic-section]{subfiles}

\begin{document}

\olfileid{cmp}{thy}{fix}
\olsection{مبرهنة النقطة الثابتة}

نعود إلى مسألة التوقف، ونتخذ مؤقتًا الرمز
$\gn{\cfind{{\OLId{latin-ordinary}{x}}}({\OLId{latin-ordinary}{y}})}$ بدل $\tuple{{\OLId{latin-ordinary}{x}},{\OLId{latin-ordinary}{y}}}$، كأنه
«اسم» لقيمة $\cfind{{\OLId{latin-ordinary}{x}}}({\OLId{latin-ordinary}{y}})$. وبه نعيد
صياغة أحد براهين عدم قابلية مشكلة التوقف للقرار.

فنسأل: أتوجد دالة قابلة للحساب ${\OLId{latin-ordinary}{h}}$ تحقق، لكل ${\OLId{latin-ordinary}{x}}$
و${\OLId{latin-ordinary}{y}}$،
\[
{\OLId{latin-ordinary}{h}}(\gn{\cfind{{\OLId{latin-ordinary}{x}}}({\OLId{latin-ordinary}{y}})}) =
\begin{cases}
{\OLMathNumeral{1}} & \text{إذا كانت $\cfind{{\OLId{latin-ordinary}{x}}}({\OLId{latin-ordinary}{y}}) \fdefined$} \\
{\OLMathNumeral{0}} & \text{خلاف ذلك.}
\end{cases}
\]

والجواب بالنفي؛ إذ لو وجدت لكانت الدالة الجزئية
\[
{\OLId{latin-ordinary}{g}}({\OLId{latin-ordinary}{x}}) \simeq
\begin{cases}
{\OLMathNumeral{0}} & \text{إذا كان ${\OLId{latin-ordinary}{h}}(\gn{\cfind{{\OLId{latin-ordinary}{x}}}({\OLId{latin-ordinary}{x}})}) = {\OLMathNumeral{0}}$} \\
\text{غير معرّفة} & \text{خلاف ذلك}
\end{cases}
\]
قابلة للحساب، ومن ثم لكان لها فهرس ما~${\OLId{latin-ordinary}{e}}$. لكن عندئذ يكون
\[
\cfind{{\OLId{latin-ordinary}{e}}}({\OLId{latin-ordinary}{e}}) \simeq
\begin{cases}
{\OLMathNumeral{0}} & \text{إذا كان ${\OLId{latin-ordinary}{h}}(\gn{\cfind{{\OLId{latin-ordinary}{e}}}({\OLId{latin-ordinary}{e}})}) = {\OLMathNumeral{0}}$} \\
\text{غير معرّفة} & \text{خلاف ذلك،}
\end{cases}
\]
فيلزم أن $\cfind{{\OLId{latin-ordinary}{e}}}({\OLId{latin-ordinary}{e}})$ معرّفة إذا وفقط إذا كانت غير معرّفة، وهو محال.

فتأمل المعادلة التي وردت فيها $\cfind{{\OLId{latin-ordinary}{e}}}$: فيها ضرب من الإحالة
الذاتية، إذ جعلنا قيمة $\cfind{{\OLId{latin-ordinary}{e}}}({\OLId{latin-ordinary}{e}})$ تعتمد، على وجه مخصوص، على
$\gn{\cfind{{\OLId{latin-ordinary}{e}}}({\OLId{latin-ordinary}{e}})}$. ومبرهنة النقطة الثابتة تثبت \emph{إمكان}
هذا الصنيع على العموم، لا في مقام البرهان بالخلف وحده.

وفي \olref{lem:fixed-equiv} عبارتان متكافئتان لمبرهنة النقطة
الثابتة. وصحيح أن صدقهما معًا يكفي، من الوجهة المنطقية، لتكافئهما؛
غير أن المراد هنا استنباط كل واحدة منهما من الأخرى مباشرة.
فبهذا تصلحان عبارتين بديلتين عن مبرهنة واحدة.

\begin{lem}
\ollabel{lem:fixed-equiv}
العبارتان الآتيتان متكافئتان:
\begin{enumerate}
\item لكل دالة جزئية قابلة للحساب ${\OLId{latin-ordinary}{g}}({\OLId{latin-ordinary}{x}},{\OLId{latin-ordinary}{y}})$ فهرس~${\OLId{latin-ordinary}{e}}$ بحيث، لكل~${\OLId{latin-ordinary}{y}}$،
\[
\cfind{{\OLId{latin-ordinary}{e}}}({\OLId{latin-ordinary}{y}}) \simeq {\OLId{latin-ordinary}{g}}({\OLId{latin-ordinary}{e}},{\OLId{latin-ordinary}{y}}).
\]
\item لكل دالة قابلة للحساب~${\OLId{latin-ordinary}{f}}({\OLId{latin-ordinary}{x}})$ فهرس~${\OLId{latin-ordinary}{e}}$ بحيث، لكل~${\OLId{latin-ordinary}{y}}$،
\[
\cfind{{\OLId{latin-ordinary}{e}}}({\OLId{latin-ordinary}{y}}) \simeq \cfind{{\OLId{latin-ordinary}{f}}({\OLId{latin-ordinary}{e}})}({\OLId{latin-ordinary}{y}}).
\]
\end{enumerate}
\end{lem}

\begin{proof}
$({\OLMathNumeral{1}}) \Rightarrow ({\OLMathNumeral{2}})$: إذا أخذنا ${\OLId{latin-ordinary}{f}}$، عرّفنا ${\OLId{latin-ordinary}{g}}$ بوضع
${\OLId{latin-ordinary}{g}}({\OLId{latin-ordinary}{x}},{\OLId{latin-ordinary}{y}}) \simeq \fn{Un}({\OLId{latin-ordinary}{f}}({\OLId{latin-ordinary}{x}}),{\OLId{latin-ordinary}{y}})$. ثم استعملنا (\OLClassicalDigits{1})، فحصل فهرس~${\OLId{latin-ordinary}{e}}$
يحقق، لكل~${\OLId{latin-ordinary}{y}}$،
\begin{align*}
\cfind{{\OLId{latin-ordinary}{e}}}({\OLId{latin-ordinary}{y}}) & = \fn{Un}({\OLId{latin-ordinary}{f}}({\OLId{latin-ordinary}{e}}),{\OLId{latin-ordinary}{y}}) \\
& = \cfind{{\OLId{latin-ordinary}{f}}({\OLId{latin-ordinary}{e}})}({\OLId{latin-ordinary}{y}}).
\end{align*}

$({\OLMathNumeral{2}}) \Rightarrow ({\OLMathNumeral{1}})$: إذا أخذنا ${\OLId{latin-ordinary}{g}}$، أعطتنا مبرهنة ${\OLId{latin-ordinary}{s}}$-${\OLId{latin-ordinary}{m}}$-${\OLId{latin-ordinary}{n}}$
دالة ${\OLId{latin-ordinary}{f}}$ تحقق، لكل ${\OLId{latin-ordinary}{x}}$ و~${\OLId{latin-ordinary}{y}}$،
$\cfind{{\OLId{latin-ordinary}{f}}({\OLId{latin-ordinary}{x}})}({\OLId{latin-ordinary}{y}}) \simeq {\OLId{latin-ordinary}{g}}({\OLId{latin-ordinary}{x}},{\OLId{latin-ordinary}{y}})$. ونستعمل (\OLClassicalDigits{2})، فنأخذ فهرسًا~${\OLId{latin-ordinary}{e}}$
بحيث
\begin{align*}
\cfind{{\OLId{latin-ordinary}{e}}}({\OLId{latin-ordinary}{y}}) & = \cfind{{\OLId{latin-ordinary}{f}}({\OLId{latin-ordinary}{e}})}({\OLId{latin-ordinary}{y}}) \\
& = {\OLId{latin-ordinary}{g}}({\OLId{latin-ordinary}{e}},{\OLId{latin-ordinary}{y}}).
\end{align*}
وبذلك يتم البرهان.
\end{proof}

\begin{explain}
وقبل إثبات (\OLClassicalDigits{1})، وما يتبعها من (\OLClassicalDigits{2})، تأمل ما في معناها من غرابة.
فلتكن ${\OLId{latin-ordinary}{e}}$ برنامجًا؛ تقول (\OLClassicalDigits{1}) إنه مهما أعطيتنا
دالة جزئية قابلة للحساب ${\OLId{latin-ordinary}{g}}({\OLId{latin-ordinary}{x}},{\OLId{latin-ordinary}{y}})$، أمكن إيجاد برنامج ${\OLId{latin-ordinary}{e}}$
يحسب ${\OLId{latin-ordinary}{g}}_{\OLId{latin-ordinary}{e}}({\OLId{latin-ordinary}{y}}) \simeq {\OLId{latin-ordinary}{g}}({\OLId{latin-ordinary}{e}},{\OLId{latin-ordinary}{y}})$. فكأننا نصنع برنامجًا يحسب
دالة ترجع إلى ذلك البرنامج نفسه.
\end{explain}

\begin{thm}
العبارتان الواردتان في \olref{lem:fixed-equiv} صادقتان. وعلى وجه
التحديد، لكل دالة جزئية قابلة للحساب ${\OLId{latin-ordinary}{g}}({\OLId{latin-ordinary}{x}},{\OLId{latin-ordinary}{y}})$ فهرس~${\OLId{latin-ordinary}{e}}$ بحيث،
لكل~${\OLId{latin-ordinary}{y}}$،
\[
\cfind{{\OLId{latin-ordinary}{e}}}({\OLId{latin-ordinary}{y}}) \simeq {\OLId{latin-ordinary}{g}}({\OLId{latin-ordinary}{e}},{\OLId{latin-ordinary}{y}}).
\]
\end{thm}

\begin{proof}
قد تضمنت مناقشة التوقف السابقة أصول هذا البرهان. فلنجعل
$\fn{diag}({\OLId{latin-ordinary}{x}})$ دالة قابلة للحساب تعيد، لكل ${\OLId{latin-ordinary}{x}}$، فهرسًا للدالة
${\OLId{latin-ordinary}{f}}_{\OLId{latin-ordinary}{x}}({\OLId{latin-ordinary}{y}}) \simeq \cfind{{\OLId{latin-ordinary}{x}}}[{\OLMathNumeral{2}}]({\OLId{latin-ordinary}{x}},{\OLId{latin-ordinary}{y}})$، أي
\[
\cfind{\fn{diag}({\OLId{latin-ordinary}{x}})}({\OLId{latin-ordinary}{y}}) \simeq \cfind{{\OLId{latin-ordinary}{x}}}[{\OLMathNumeral{2}}]({\OLId{latin-ordinary}{x}},{\OLId{latin-ordinary}{y}}).
\]
والدالة $\fn{diag}$ تحول برنامج دالة ثنائية إلى برنامج
دالة أحادية، بأن تثبت البرنامج الأول حجةً أولى. ولتعريف
$\fn{diag}$ تعريفًا صوريًا، نبدأ بتعريف ${\OLId{latin-ordinary}{s}}$ بوضع
\[
{\OLId{latin-ordinary}{s}}({\OLId{latin-ordinary}{x}},{\OLId{latin-ordinary}{y}}) \simeq \fn{Un}^{\OLMathNumeral{2}}({\OLId{latin-ordinary}{x}},{\OLId{latin-ordinary}{x}},{\OLId{latin-ordinary}{y}}),
\]
حيث $\fn{Un}^{\OLMathNumeral{2}}$ دالة ثلاثية شاملة للدوال الجزئية الثنائية القابلة
للحساب. ومن مبرهنة ${\OLId{latin-ordinary}{s}}$-${\OLId{latin-ordinary}{m}}$-${\OLId{latin-ordinary}{n}}$ نستخرج دالة عودية بدائية
$\fn{diag}$ تحقق
\[
\cfind{\fn{diag}({\OLId{latin-ordinary}{x}})}({\OLId{latin-ordinary}{y}}) \simeq {\OLId{latin-ordinary}{s}}({\OLId{latin-ordinary}{x}},{\OLId{latin-ordinary}{y}}).
\]

ثم نعرّف ${\OLId{latin-ordinary}{l}}$ بالمعادلة
\[
{\OLId{latin-ordinary}{l}}({\OLId{latin-ordinary}{x}},{\OLId{latin-ordinary}{y}}) \simeq {\OLId{latin-ordinary}{g}}(\fn{diag}({\OLId{latin-ordinary}{x}}),{\OLId{latin-ordinary}{y}}).
\]
ونجعل $\gn{{\OLId{latin-ordinary}{l}}}$ فهرسًا لـ${\OLId{latin-ordinary}{l}}$، ونضع أخيرًا
${\OLId{latin-ordinary}{e}}=\fn{diag}(\gn{{\OLId{latin-ordinary}{l}}})$. عندئذ، لكل ${\OLId{latin-ordinary}{y}}$، لدينا
\begin{align*}
\cfind{{\OLId{latin-ordinary}{e}}}({\OLId{latin-ordinary}{y}}) & \simeq \cfind{\fn{diag}(\gn{{\OLId{latin-ordinary}{l}}})}({\OLId{latin-ordinary}{y}}) \\
& \simeq \cfind{\gn{{\OLId{latin-ordinary}{l}}}}[{\OLMathNumeral{2}}](\gn{{\OLId{latin-ordinary}{l}}}, {\OLId{latin-ordinary}{y}}) \\
& \simeq {\OLId{latin-ordinary}{l}}(\gn{{\OLId{latin-ordinary}{l}}}, {\OLId{latin-ordinary}{y}}) \\
& \simeq {\OLId{latin-ordinary}{g}}(\fn{diag}(\gn{{\OLId{latin-ordinary}{l}}}),{\OLId{latin-ordinary}{y}}) \\
& \simeq {\OLId{latin-ordinary}{g}}({\OLId{latin-ordinary}{e}}, {\OLId{latin-ordinary}{y}}),
\end{align*}
كما هو مطلوب.
\end{proof}

\begin{explain}
فما وجه هذا البناء؟ لنفرض أنك تريد برنامجًا يطبع نفسه،
وأن لغة البرمجة التي تستعملها ذات مكتبة واسعة، غير مألوفة، من دوال
السلاسل. ومن جملتها دالة $\fn{diag}$، عملها
أن تتلقى السلسلة~${\OLId{latin-ordinary}{s}}$، فتعيّن $\fn{diag}$ كل موضع يرد فيه
الرمز `x' في~${\OLId{latin-ordinary}{s}}$ وتضع مكانه نسخة مقتبسة من السلسلة الأصلية. فإذا
أدخلنا، مثلًا، السلسلة
\begin{quote}
\begin{verbatim}
hello x world
\end{verbatim}
\end{quote}
دخلًا، أعادت الدالة
\begin{quote}
\begin{verbatim}
hello 'hello x world' world
\end{verbatim}
\end{quote}
خرجًا. فيسهل حينئذ إنشاء البرنامج المطلوب؛ وتحقق أن
\begin{quote}
\begin{verbatim}
print(diag('print(diag(x))'))
\end{verbatim}
\end{quote}
يؤدي الغرض. وتجري الفكرة ذاتها في لغات مألوفة، مثل C++ وJava،
وإن زاد تفصيل تنفيذها.

ومن هنا نصل إلى برهان النقطة الثابتة بخطوتين. لنفرض أن
لدالة الطباعة صورة معدلة $\fn{print}({\OLId{latin-ordinary}{x}},{\OLId{latin-ordinary}{y}})$ تتلقى سلسلة ${\OLId{latin-ordinary}{x}}$ وحجة
عددية ${\OLId{latin-ordinary}{y}}$، وتطبع السلسلة ${\OLId{latin-ordinary}{x}}$ متكررة ${\OLId{latin-ordinary}{y}}$ مرة.
فيكون «البرنامج»
\begin{quote}
\begin{verbatim}
getinput(y);
print(diag('getinput(y); print(diag(x), y)'), y)
\end{verbatim}
\end{quote}
برنامجًا يطبع نفسه ${\OLId{latin-ordinary}{y}}$ مرة عند الدخل ${\OLId{latin-ordinary}{y}}$. فإذا أقمنا مقام الهيكل
$\fn{getinput}$---$\fn{print}$---$\fn{diag}$ دالة اعتباطية
${\OLId{latin-ordinary}{g}}({\OLId{latin-ordinary}{x}},{\OLId{latin-ordinary}{y}})$ حصل
\begin{quote}
\begin{verbatim}
g(diag('g(diag(x), y)'), y)
\end{verbatim}
\end{quote}
فهذا برنامج يجري، عند الدخل ${\OLId{latin-ordinary}{y}}$، الدالة ${\OLId{latin-ordinary}{g}}$ على البرنامج نفسه مع
${\OLId{latin-ordinary}{y}}$. ومتى فسرنا «الاقتباس» بأنه «استعمال فهرس لـ»، رجعنا إلى
البرهان المتقدم.

ولا يضر الآن أن يبدو لك البرهان حيلة صورية أو شيئًا من
السحر؛ وإنما المطلوب أن تقدر على إعادة تفاصيل الحجة. وعند
إثبات مبرهنات عدم الاكتمال، و«مبرهنة النقطة الثابتة» المتصلة بها،
ننظر في وجوه أخرى تفسر نجاح هذا البناء.
\end{explain}

\begin{tagblock}{lambda}
\begin{digress}
وبالفكرة نفسها نحصل على مؤلّف «نقطة ثابتة». فإذا
أخذنا حد لامبدا ${\OLId{latin-ordinary}{g}}$، وطلبنا حدًا ${\OLId{latin-ordinary}{k}}$ يكون فيه ${\OLId{latin-ordinary}{k}}$ مكافئًا
$\beta$ لـ$gk$، عرّفنا الحدين
\[
\fn{diag}({\OLId{latin-ordinary}{x}}) = xx
\]
و
\[
{\OLId{latin-ordinary}{l}}({\OLId{latin-ordinary}{x}}) = {\OLId{latin-ordinary}{g}}(\fn{diag}({\OLId{latin-ordinary}{x}}))
\]
بحسب اصطلاحات الترميز المتقدمة؛ فيكون ${\OLId{latin-ordinary}{l}}$ الحد
$\lambd[{\OLId{latin-ordinary}{x}}][{\OLId{latin-ordinary}{g}}(xx)]$. ولنجعل ${\OLId{latin-ordinary}{k}}$ مساويًا للحد $ll$؛ فيحصل
\begin{align*}
{\OLId{latin-ordinary}{k}} & = (\lambd[{\OLId{latin-ordinary}{x}}][{\OLId{latin-ordinary}{g}}(xx)])(\lambd[{\OLId{latin-ordinary}{x}}][{\OLId{latin-ordinary}{g}}(xx)]) \\
& \red  {\OLId{latin-ordinary}{g}}((\lambd[{\OLId{latin-ordinary}{x}}][{\OLId{latin-ordinary}{g}}(xx)])(\lambd[{\OLId{latin-ordinary}{x}}][{\OLId{latin-ordinary}{g}}(xx)])) \\
& = gk.
\end{align*}
فإذا أخذنا
\[
{\OLId{latin-looped}{Y}} = \lambd[{\OLId{latin-ordinary}{g}}][((\lambd[{\OLId{latin-ordinary}{x}}][{\OLId{latin-ordinary}{g}}(xx)])(\lambd[{\OLId{latin-ordinary}{x}}][{\OLId{latin-ordinary}{g}}(xx)]))]
\]
اختزل $Yg$ و${\OLId{latin-ordinary}{g}}(Yg)$ إلى حد مشترك؛ ومن ثم
$Yg \equiv_\beta {\OLId{latin-ordinary}{g}}(Yg)$. ويعرف هذا باسم «مؤلّف كاري». أما إذا أخذنا
بدلًا منه
\[
{\OLId{latin-looped}{Y}} = (\lambd[xg][{\OLId{latin-ordinary}{g}}(xxg)])(\lambd[xg][{\OLId{latin-ordinary}{g}}(xxg)])
\]
فإن $Yg$ يختزل في الواقع إلى ${\OLId{latin-ordinary}{g}}(Yg)$، وهذا أقوى من مجرد التكافؤ. وهذه
الصورة الأخيرة لـ${\OLId{latin-looped}{Y}}$ تسمى «مؤلّف تورنغ».
\end{digress}
\end{tagblock}

\end{document}
